\subsection{环的素谱}

设 A 是一个环，X 是 A 的素理想之集。对 A 的每个子集 M，我们用 V(M) 表示包含 M 的 A 的素理想之集；显然，若 a 是 M 生成的 A 的理想，则 V(M) = V(a)；若 M 由单个元素 f 组成，则我们用 V(f) 代替 V(\{f\})，且有 V(f) = V(Af)。映射 M \(\mapsto\) V(M) 关于 A 与 X 中的包含关系是递减的。此外，下列公式成立：

\begin{enumerate}
\def\labelenumi{(\arabic{enumi})}
\tightlist
\item
\end{enumerate}

\[
\mathbf {V} (0) = \mathbf {X}, \quad \mathbf {V} (1) = \varnothing ;\tag{1}
\]

\[
\mathrm{V} \left(\bigcup_ {i \in I} \mathrm{M} _ {i}\right) = \mathrm{V} \left(\sum_ {i \in I} \mathrm{M} _ {i}\right) = \bigcap_ {i \in I} \mathrm{V} (\mathrm{M} _ {i})\tag{2}
\]

对 A 的子集的每个族 \((\mathbf{M}_i)_{i\in I}\) 成立；

\[
\mathbf {V} (\mathfrak {a} \cap \mathfrak {a} ^ {\prime}) = \mathbf {V} (\mathfrak {a a} ^ {\prime}) = \mathbf {V} (\mathfrak {a}) \cup \mathbf {V} (\mathfrak {a} ^ {\prime})\tag{3}
\]

对 A 中理想的每个对 \(\mathfrak{a},\mathfrak{a}'\)。公式 (1) 与 (2) 是显然的；另一方面，公式 (3) 表示：为使 A 的素理想 \(\mathfrak{p}\) 包含理想 \(\mathfrak{a}\) 与 \(\mathfrak{a}'\) 之一，必须且只需它包含 \(\mathfrak{a}\mathfrak{a}'\) 或包含 \(\mathfrak{a} \cap \mathfrak{a}'\)；于是它是 §1 第 1 节命题 1 的推论。(1) 中第二个公式有下述逆命题：若 \(\mathfrak{a}\) 是 \(\mathbf{A}\) 的一个理想，使 \(V(\mathfrak{a}) = \varnothing\)，则 \(\mathfrak{a} = \mathbf{A}\)，因为不存在 \(\mathbf{A}\) 的包含 \(\mathfrak{a}\) 的极大理想。最后，若 \(\mathfrak{a}\) 是 \(\mathbf{A}\) 的一个理想，而 \(r(\mathfrak{a})\) 是它的根（§2 第 6 节定义 4），则

\[
\mathbf {V} (\mathfrak {a}) = \mathbf {V} (\mathfrak {r} (\mathfrak {a}))\tag{4}
\]

这由 §2 第 6 节命题 13 的推论 1 得出。

公式 (1) 至 (3) 表明，X 的子集 V(M) 满足一个拓扑的闭集公理（《一般拓扑学》第 I 章 §1 第 4 节）。

定义 4. 设 A 是一个环。A 的素理想之集 X，带上以诸 V(M)（M 跑遍 P(A)）为闭集的拓扑，称为 A 的素谱，记作 Spec(A)。这样定义的拓扑称为 X 上的谱拓扑或 Zariski 拓扑。

显然，关系 \(\operatorname{Spec}(\mathbf{A}) = \varnothing\) 等价于 \(\mathbf{A} = \{0\}\)。

设 X 是环 A 的素谱；对每个 \(f \in A\)，用 \(X_{f}\) 表示不包含 f 的 A 的素理想之集；则 \(\mathbf{X}_{f} = \mathbf{X} - \mathbf{V}(f)\)，因而 \(X_{f}\) 是开集。由 (2)，X 的每个闭子集都是形如 \(\mathrm{V}(f)\) 的闭集之交，故诸 \(X_{f}\) 构成 X 上谱拓扑的一个基。此外，由定义立即得出

\[
\mathbf {X} _ {0} = \varnothing , \quad \mathbf {X} _ {1} = \mathbf {X},\tag{5}
\]

更一般地，对 A 的每个可逆元 f 有 \(X_{f} = X\)；

\[
\mathrm{X} _ {f g} = \mathrm{X} _ {f} \cap \mathrm{X} _ {g} \quad \text {   for   all   } f, g \text {   in   A.   }\tag{6}
\]

对 X 的每个子集 Y，用 \(\Im(Y)\) 表示属于 Y 的 A 的素理想之交。显然 \(\Im(Y)\) 是 A 的理想，且映射 \(Y \mapsto \Im(Y)\) 关于 X 与 A 中的包含关系是递减的。显然，关系

\begin{enumerate}
\def\labelenumi{(\arabic{enumi})}
\setcounter{enumi}{6}
\tightlist
\item
\end{enumerate}

\[
\Im (\varnothing) = \mathbf {A}\tag{7}
\]

\[
\Im \left(\bigcup_ {\lambda \in L} Y _ {\lambda}\right) = \bigcap_ {\lambda \in L} \Im (Y _ {\lambda})\tag{8}
\]

对族 \((\mathbf{Y}_{\lambda})_{\lambda \in \mathbf{L}}\)（\(\mathbf{X}\) 的子集组成的族）成立。此外：

命题 11. 设 A 是一个环，a 是 A 的一个理想，Y 是 X = Spec(A) 的一个子集。

\begin{enumerate}
\def\labelenumi{(\roman{enumi})}
\item
  \(\mathbf{V}(\mathfrak{a})\) 在 \(\mathbf{X}\) 中闭，而 \(\Im (\mathbf{Y})\) 是 A 的一个等于其根的理想。
\item
  \(\Im(\mathrm{V}(\mathfrak{a}))\) 是 \(\mathfrak{a}\) 的根，且 \(\mathrm{V}(\Im(\mathrm{Y}))\) 是 \(\mathrm{Y}\) 在 \(\mathrm{X}\) 中的闭包。
\item
  映射 \(\Im\) 与 \(\mathbf{V}\) 定义互逆的两个递减双射，它们建立于 \(\mathbf{X}\) 的闭子集之集与 \(\mathbf{A}\) 的等于其根的理想之集之间。
\end{enumerate}

断言 (i) 与 (ii) 的第一个断言由定义及 §2 第 6 节命题 13 的推论 1 得出。若闭集 V(M)（对某个 M ⊂ A）包含 Y，则对每个素理想 p ∈ Y 有 M ⊂ p，从而 M ⊂ \(\mathfrak{I}\)(Y)，因而 V(M) ⊃ V(\(\mathfrak{I}\)(Y))；由于 Y ⊂ V(\(\mathfrak{I}\)(Y))，V(\(\mathfrak{I}\)(Y)) 是 X 中包含 Y 的最小闭集，这就完成了 (ii) 的证明。最后，由 (ii) 得出：若 a 是等于其根的一个理想，则 \(\mathfrak{I}\)(V(a)) = a；若 Y 在 X 中闭，则 V(\(\mathfrak{I}\)(Y)) = Y。这就证明了 (iii)。

由命题 11 立即得出：若 M 是 A 的任意子集而 Y 是 X 的任意子集，则 V(M) = V(\(\Im(V(M))\)) 且 \(\Im(Y) = \Im(V(\Im(Y)))\)。

推论 1. 对每个族 \((\mathbf{Y}_{\lambda})_{\lambda \in L}\)（由 \(\mathbf{X}\) 的闭子集组成），\(\Im\left(\bigcap_{\lambda \in L} \mathbf{Y}_{\lambda}\right)\) 是诸理想 \(\Im(\mathbf{Y}_{\lambda})\) 之和的根。

由命题 11 (iii)，\(\Im\left(\bigcap_{\lambda\in L}Y_{\lambda}\right)\) 是等于其根且包含所有 \(\Im(Y_{\lambda})\) 的最小理想；于是该理想包含 \(\sum_{\lambda\in L}\Im(Y_{\lambda})\)，从而也包含 \(\sum_{\lambda\in L}\Im(Y_{\lambda})\) 的根（\(\S\) 2 第 6 节命题 13 的推论 2），由此得该推论。

推论 2. 以 \(\mathfrak{r}(\mathfrak{a})\) 表示理想 \(\mathfrak{a}\)（在 \(\mathbf{A}\) 中）的根；若 \(\mathfrak{a}\) 与 \(\mathfrak{b}\) 是 \(\mathbf{A}\) 的两个理想，则关系 \(\mathrm{V}(\mathfrak{a}) \subset \mathrm{V}(\mathfrak{b})\) 等价于 \(\mathfrak{b} \subset \mathfrak{r}(\mathfrak{a})\) 与 \(\mathfrak{r}(\mathfrak{b}) \subset \mathfrak{r}(\mathfrak{a})\)。

关系 \(\mathfrak{b} \subset \mathfrak{r}(\mathfrak{a})\) 与 \(r(\mathfrak{b}) \subset \mathfrak{r}(\mathfrak{a})\) 等价是直接的，且由于 \(\mathrm{V}(\mathfrak{a}) = \mathrm{V}(\mathfrak{r}(\mathfrak{a}))\)，该推论立即由命题 11 (iii) 得出。

推论 3. 设 \((f_{\lambda})_{\lambda \in L}\) 是 A 的元素的一个族。为使元素 \(g \in \mathbf{A}\) 满足 \(\mathbf{X}_g \subset \bigcup_{\lambda \in L} \mathbf{X}_{f_\lambda}\)，必须且只需存在整数 \(n > 0\) 使 \(g^n\) 属于 \(f_\lambda\) 生成的理想。

关系 \(X_{g} \subset \bigcup_{\lambda \in L} X_{f_{\lambda}}\) 等价于 \(\mathrm{V}(g) \supset \bigcap_{\lambda \in L} \mathrm{V}(f_{\lambda})\)，只需应用推论 2 即可。

推论 4. 为使 \(f, g\)（\(A\) 的两个元素）满足 \(X_f = X_g\)，必须且只需存在两个整数 \(m > 0, n > 0\) 使 \(f^m \in Ag\) 且 \(g^n \in Af\)。

* 推论 5. 为使 \(f \in \mathbf{A}\) 满足 \(\mathbf{X}_f = \varnothing\)，必须且只需 \(f\) 是幂零的。

这由推论 4 立即得出。

推论 6. 由单个点 \(\mathfrak{p} \in \mathbf{X} = \operatorname{Spec}(\mathbf{A})\) 组成的集的闭包是集 \(\mathbf{V}(\mathfrak{p})\)（包含 \(\mathfrak{p}\) 的素理想之集）。为使集 \(\{\mathfrak{p}\}\) 在 \(\mathbf{X}\) 中闭（或者，按我们也将采用的一种滥用说法，为使 \(\mathfrak{p}\) 是 \(\mathbf{X}\) 的一个闭点），必须且只需 \(\mathfrak{p}\) 是极大的。

推论 7. 若 A 是 Noether 环，则 X = Spec(A) 是 Noether 空间。

命题 12. 对每个 \(f \in \mathbf{A}\)，\(\mathbf{X}_f\)（\(\mathbf{X} = \operatorname{Spec}(\mathbf{A})\) 中的开集）是拟紧的；特别地，空间 \(\mathbf{X}\) 是拟紧的。

由于诸 \(X_{g}\) 构成该拓扑的基，只需证明：若 \((g_{\lambda})_{\lambda\in L}\) 是 A 的元素的族，使 \(X_{f}\subset\bigcup_{\lambda\in L}X_{g_{\lambda}}\)，则存在有限子族 \((g_{\lambda})_{\lambda\in H}\) 使 \(X_{f}\subset\bigcup_{\lambda\in H}X_{g_{\lambda}}\)。但关系 \(X_{f}\subset\bigcup_{\lambda\in L}X_{g_{\lambda}}\) 表示存在整数 \(n > 0\) 与有限子族 \((g_{\lambda})_{\lambda\in L}\) 使 \(f^{n}\) 属于该子族生成的理想（命题 11 的推论 3）；由此得该命题。

命题 13. 设 A、A' 是两个环，X = Spec(A)，X' = Spec(A')，h 是从 A 到 A' 的一个同态；映射 \(^a h\)：\(p' \mapsto \overline{h}^{-1}(p')\)（从 X' 到 X）是连续的。

对 \(\mathbf{M} \subset \mathbf{A}\)，集合 \((^a h)^{-1}(\mathbf{V}(\mathbf{M}))\) 是素理想 \(\mathfrak{p}'\)（属于 \(\mathbf{A}'\)，且满足 \(\mathbf{M} \subset \overline{h}^{-1}(\mathfrak{p}')\)）所成之集，这等价于 \(h(\mathbf{M}) \subset \mathfrak{p}'\)；该集合于是等于 \(\mathbf{V}(h(\mathbf{M}))\)，因而是闭的。

我们把 \(^{a}h\) 称为与同态 h 相伴的映射。

注. 若 \(h\) 是满射且 \(\mathfrak{a}\) 是它的核，则由谱拓扑的定义得出，\(^a h\) : 是从 \(X'\) 到闭子空间 \(V(\mathfrak{a})\)（\(X\) 的）上的一个同胚；为使 \(\mathfrak{p}'\)（\(A'\) 的素理想）包含 \(\mathfrak{b}'\)（\(A'\) 的理想），必须且只需 \(\overline{h}^{-1}(\mathfrak{p}')\) 包含 \(\overline{h}^{-1}(\mathfrak{b}')\)；首先，\(^a h\) 是单射（取 \(\mathfrak{b}'\) 为素理想即见）；此外，对每个理想 \(\mathfrak{b}'\)（属于 \(A'\)），\(^a h\) 下 \(V(\mathfrak{b}')\) 的像是 \(V(\overline{h}^{-1}(\mathfrak{b}'))\)，由此得我们的断言，因为形如 \(\overline{h}^{-1}(\mathfrak{b}')\) 的理想都是 \(A\) 的包含 \(\mathfrak{a}\) 的理想。

推论. 设 S 是 A 的一个乘性子集，\(A' = S^{-1}A\)，h 是典范同态 \(i_{A}^{S}\)；则 \(^{a}h\) 是从 \(X' = \operatorname{Spec}(A')\) 到 \(X = \operatorname{Spec}(A)\) 的子空间（由不与 S 相交的 A 的素理想组成）上的一个同胚。

设 \(f' = f / s\)，其中 \(f \in \mathbf{A}\)，\(s \in \mathbf{S}\)；则 \(\mathbf{X}_{f'}' = \mathbf{X}_{f/1}'\)，因为 \(s / 1\) 在 \(\mathbf{A}'\) 中可逆。我们已经知道 \(^a h\) 是单射，且对一切 \(p' \in \mathbf{X}'\)，关系 \(f / 1 \in p'\) 与 \(f \in \bar{h}^{-1}(p') = ^a h(p')\) 等价，因而条件 \(p' \in \mathbf{X}_{f/1}'\) 与 \(^a h(p') \in \mathbf{X}_f\) 等价；这表明 \(^a h(\mathbf{X}_{f'})\) 等于 \(\mathbf{X}_f \cap ^a h(\mathbf{X}')\)，由此得第一个断言，因为诸 \(\mathbf{X}_f\)（相应地，\(\mathbf{X}_{f'}'\)）构成 \(\mathbf{X}\)（相应地，\(\mathbf{X}'\)）的拓扑的基。第二个断言由 §2 第 5 节命题 11 (ii) 得出。

命题 14. 设 A 是一个环。为使 X = Spec(A) 的子集 Y 不可约，必须且只需理想 \(\Im(Y)\) 是素的。

记 \(\mathfrak{p} = \Im(\mathbf{Y})\)。我们注意，对元素 \(f \in \mathbf{A}\)，关系 \(f \in \mathfrak{p}\) 等价于 \(\mathbf{Y} \subset \mathbf{V}(f)\)。设 \(\mathbf{Y}\) 不可约，并设 \(f, g\) 是 \(\mathbf{A}\) 中使 \(fg \in \mathfrak{p}\) 的元素。则

\[
\mathbf {Y} \subset \mathbf {V} (f g) = \mathbf {V} (f) \cup \mathbf {V} (g);
\]

由于 Y 不可约且 V(f) 与 V(g) 是闭的，故 Y ⊂ V(f) 或 Y ⊂ V(g)，因而 f ∈ p 或 g ∈ p，这证明 p 是素的。

现在设 p 是素的；则 \(\overline{\mathbf{Y}} = \mathbf{V}(\mathfrak{p})\)（命题 11 (ii)），且由于 p 是素的，\(\mathfrak{p} = \Im(\{\mathfrak{p}\})\)，从而 \(\overline{\mathbf{Y}} = \mathbf{V}(\Im(\{\mathfrak{p}\})) = \{\mathfrak{p}\}\)（命题 11 (ii)）。由于由单个点组成的集是不可约的，故 Y 不可约（第 1 节命题 2）。

推论 1. 为使环 A 满足 \(\mathbf{X} = \operatorname{Spec}(\mathbf{A})\) 不可约，必须且只需 A 被其幂零根 \(\mathfrak{N}\) 除所得的商是整环。

由命题 11 (i)，\(\Im(\mathbf{X})\) 是理想 (0) 的根，即 \(\mathfrak{N}\)。

推论 2. 映射 \(\mathfrak{p} \mapsto \mathrm{V}(\mathfrak{p})\) 是从 \(\mathrm{X} = \operatorname{Spec}(\mathrm{A})\) 到 \(\mathrm{X}\) 的不可约闭子集之集上的双射；特别地，闭子集 \(\mathrm{Y}\)（\(\mathrm{X}\) 中的）的不可约分支是诸集 \(\mathrm{V}(\mathfrak{p})\)，其中 \(\mathfrak{p}\) 跑遍 \(\mathrm{A}\) 的包含 \(\Im(\mathrm{Y})\) 的素理想之集的极小元。

由于 \(\Im (\mathbf{V}(\mathfrak{p})) = \mathfrak{p}\) 对 \(\mathfrak{p}\)（\(\mathbf{A}\) 的每个素理想）成立，且 \(\mathbf{Y} = \mathbf{V}(\Im (\mathbf{Y}))\) 对 \(\mathbf{Y}\)（\(\mathbf{X}\) 的每个闭子集）成立，第一个断言由命题 14 得出；另一方面，为使 \(\mathbf{Y} \supset \mathbf{V}(\mathfrak{p})\)，必须且只需

\[
\mathfrak {p} = \Im (\mathbf {V} (\mathfrak {p})) \supset \Im (\mathbf {Y})
\]

（命题 11），由此得第二个断言。

推论 3. Noether 环 A 的极小素理想之集是有限的。

X = Spec(A) 只有有限多个不可约分支（命题 11 的推论 7 与第 2 节命题 10），该推论由上面的推论 2 得出。

命题 15. 设 A 是一个环，I 是一个有限集，E 是正交族 \((e_i)_{i \in I}\)（由 A 的满足 \(e_i \neq 0\) 的幂等元组成，且 \(\sum_{i \in I} e_i = 1\)）全体之集。对每个 \((e_i)_{i \in I} \in E\)，我们置 \(\overline{\omega}((e_i)_{i \in I}) = (V(A(1 - e_i)))_{i \in I}, \sigma((e_i)_{i \in I}) = (Ae_i)_{i \in I}\)。则 \(\overline{\omega}\) 是从 E 到分划 \((U_i)_{i \in I}\)（把 \(X = \text{Spec}(A)\) 分成开集）组成的集 P 上的双射，而 \(\sigma\) 是从 E 到族 \((a_i)_{i \in I}\)（由 A 的理想 \(\neq 0\) 组成，使 A 是诸 \(a_i\) 的直和）组成的集 S 上的双射。

设 \((e_i)_{i\in I}\) 是 \(E\) 的一个元素，并置 \(\mathbf{Y}_i = \mathbf{V}(\mathbf{A}(1 - e_i))\)；若 \(i\neq j\)，则 \(1 = 1 - e_i + e_i(1 - e_j)\in A(1 - e_i) + A(1 - e_j)\)，从而 \(\mathbf{Y}_i\cap \mathbf{Y}_j = \varnothing\)（公式 (1) 与 (2)）。另一方面，

\[
\bigcup_ {i \in I} Y _ {i} = V \left(\prod_ {i \in I} A (1 - e _ {i})\right) \quad (\text { formula   (3) });
\]

按假设 \(\prod_{i\in I}(1 - e_i) = 1 - \sum_{i\in I}e_i = 0\)，从而 \(\bigcup_{i\in I}\mathbf{Y}_i = \mathbf{X}\)（公式 (1)）。由于诸 \(\mathbf{Y}_i\) 是闭的，它们也是开的，从而 \(\overline{\omega} (\mathrm{E})\subset \mathrm{P}\)。另外，显然 \(\mathbf{A} = \sum_{i\in I}\mathbf{A}e_i\)；若 \(0 = \sum_{i\in I}a_{i}e_{i}\)，其中 \(a_{i}\in \mathbf{A}\)，则用 \(e_i\) 乘之得 \(0 = a_i e_i^2 = a_i e_i\) 对一切 \(i\) 成立；这证明 \(\sigma (\mathrm{E})\subset \mathrm{S}\)。

引理 2. 若 \(e, f\) 是 \(A\) 的两个幂等元，使 \(Ae\) 与 \(Af\) 有相同的根，则 \(e = f\)。

按假设存在整数 \(m \geqslant 0\)、\(n \geqslant 0\) 使 \(e = e^{m} \in Af\) 且 \(f = f^{n} \in Ae\)；设 x、y 是 A 中使 e = xf、f = ye 的元素；则 ef = \(xf^{2} = xf = e\) 且类似地 ef = \(ye^{2} = ye = f\)，从而 e = f。

引理 2 与命题 11 的推论 2 表明映射 \(\overline{\omega}\) 与 \(\sigma\) 是单射。

我们证明 \(\sigma\) 是满射。若 \((a_{i})_{i\in I}\) 是 \(S\) 的一个元素，则存在元素 \(e_i\in a_i\) 使 \(1 = \sum_{i\in I}e_i\)；若 \(i\neq j\)，则 \(e_i e_j\in a_i\cap a_j = \{0\}\)，从而

\[
e _ {i} = \sum_ {j \in I} e _ {i} e _ {j} = e _ {i} ^ {2};
\]

最后，\(\mathbf{A}e_{i}\subset \mathfrak{a}_{i}\) 对一切 \(i\in \mathbf{I}\) 成立，且 \(\sum_{i\in \mathbf{I}}\mathrm{A}e_i = \mathrm{A},\)，从而 \(\mathrm{A}e_{i} = \mathfrak{a}_{i}\)。

剩下要证 \(\overline{\omega}\) 是满射。设 \((\mathbf{U}_i)_{i\in I}\) 是 \(P\) 的一个元素，并置 \(Z_{i} = \complement U_{i} = \bigcup_{j\neq i}U_{j}\)；由于 \(U_{i}\) 与 \(Z_{i}\) 是闭的，存在理想 \(\mathfrak{a}_i,\mathfrak{b}_i\)（\(A\) 的）使 \(\mathbf{U}_i = \mathbf{V}(\mathfrak{a}_i)\)、\(Z_{i} = \mathbf{V}(\mathfrak{b}_{i})\)。我们现在证明还可以进一步假设 \(\mathfrak{a}_i\cap \mathfrak{b}_i = 0\)。现在 \(\mathbf{U}_i\cap Z_i = \varnothing\)，从而 \(\mathfrak{a}_i + \mathfrak{b}_i = A\)；设 \(a_i\in \mathfrak{a}_i\)、\(b_{i}\in \mathfrak{b}_{i}\) 使 \(a_{i} + b_{i} = 1\)。则 \(\mathbf{X} = \mathbf{U}_i\cup Z_i = \mathbf{V}(\mathfrak{a}_i\mathfrak{b}_i)\)（公式 (3)）；因此 \(\mathfrak{a}_i\mathfrak{b}_i\) 的每个元素都是幂零的（命题 11 的推论 2）；设 \(p\) 是使 \(a_i^p b_i^p = 0\) 的整数。则 \(\mathbf{U}_i\subset \mathbf{V}(\mathbf{A}a_i) = \mathbf{V}(\mathbf{A}a_i^p)\)，

\[
\mathbf {Z} _ {i} \subset \mathbf {V} (\mathbf {A} b _ {i}) = \mathbf {V} (\mathbf {A} b _ {i} ^ {p})
\]

且 \(\mathrm{V}(\mathbf{A}a_i)\cap \mathrm{V}(\mathbf{A}b_i) = \mathrm{V}(\mathbf{A}a_i + \mathbf{A}b_i) = \varnothing\)，因而 \(\mathbf{U}_i = \mathrm{V}(\mathbf{A}a_i^p)\) 且 \(\mathbf{Z}_i = \mathrm{V}(\mathbf{A}b_i^p)\)，把 \(a_i\) 换成 \(\mathbf{A}a_i^p\)、\(b_i\) 换成 \(\mathbf{A}b_i^p\) 即可确立我们的断言。这样选定理想 \(a_i\) 与 \(b_i\) 后，由 \(\sigma\) 是双射这一事实得出，存在两个幂等元 \(f_i\in a_i\)、\(e_i\in b_i\) 使 \(1 = e_i + f_i\)、\(e_i f_i = 0\)、\(a_i = Af_i\)、\(b_i = Ae_i\)。若 \(i\neq j\)，则 \(\mathbf{X} = \mathbf{Z}_i\cup \mathbf{Z}_j = \mathrm{V}(\mathbf{A}e_i e_j)\)，且由于 \(e_i e_j\) 是幂等的，引理 2 表明 \(e_i e_j = 0\)。最后 \(e = \sum_{i\in I}e_i\) 是幂等的，且 \(e_i\in \mathbf{A}e\) 对一切 \(i\in I\) 成立，从而 \(\mathrm{V}(\mathbf{A}e)\subset \mathbf{Z}_i\) 对一切 \(i\) 成立；由此得出

\[
\mathbf {V} (\mathbf {A} e) = \varnothing = \mathbf {V} (\mathbf {A}. 1)
\]

且引理 2 还表明 \(e = 1\)。

推论 1. 设 A 是一个环，r 是 A 的一个诣零理想，h: A → A/r 是典范同态。对 A/r 的幂等元的每个有限正交族 \((e_{i}^{\prime})_{i\in I}\)（满足 \(\sum_{i\in I} e_{i}^{\prime}=1\)），存在 A 的幂等元的有限正交族 \((e_{i})_{i\in I}\)，使 \(\sum_{i\in I} e_{i}=1\)，且 \(h(e_{i})=e_{i}^{\prime}\) 对一切 \(i\in I\) 成立。

我们记 \(\mathbf{A}' = \mathbf{A} / \mathfrak{r}\)。我们知道（命题 13 后的注）

\[
^ {a} h \colon \operatorname{Spec} (\mathrm{A} ^ {\prime}) \rightarrow \operatorname{Spec} (\mathrm{A})
\]

是同胚的，因为按假设 A 的每个素理想包含 \(\mathfrak{r}\)。命题 15 表明，A 中存在一个由幂等元组成的有限正交族 \((e_i)_{i\in I}\)，使 \(\sum_{i\in I}e_i = 1\)，且在 \(^ah\) 之下，\(\mathrm{V(A'(1 - e_i'))}\) 的像是 \(\mathrm{V(A(1 - e_i))}\)。但显然 \(\mathrm{V(A(1 - e_i))}\) 也是在 \(^ah\) 之下

\[
\mathbf {V} (\mathbf {A} ^ {\prime} (1 - h (e _ {i})))  ;
\]

由于 \(1 - e_i'\) 与 \(1 - h(e_i)\) 是幂等的，引理 2 表明 \(e_i' = h(e_i)\)，由此得该推论。

推论 2. 为使素谱 \(\mathbf{X} = \operatorname{Spec}(\mathbf{A})\)（环 \(\mathbf{A}\) 的）是连通的，必须且只需 \(\mathbf{A}\) 不含 0 与 1 以外的幂等元。

说 X 不连通，意味着 X 中存在一个既开又闭且异于 \(\varnothing\) 与 X 的集合。

